\documentclass[11pt]{article}
\usepackage{amsmath,amssymb,amsthm}
\newtheorem{lemma}{Lemma}
\newtheorem{proposition}[lemma]{Proposition}
\newcommand{\gen}[1]{\langle #1\rangle}
\begin{document}
\title{A supplied invariant subset for witness 290}
\author{Written application for the finite intake}
\date{}
\maketitle

Retain the exact original degree-$3510$ action of the named inner simple
group $G=\mathrm{Fi}_{22}$ and its normal embedding in the acquired ambient
group $A$. The original model, faithful action, ordered witness, membership,
$H=B$ interpretation and two-fibre transfer are explicit imported premises.
A point list or an owner integer does not establish them. The present text
is an unissued application; it supplies no Root adoption or ledger credit.

\begin{lemma}
Let $a,b,c$ be involutions on a finite set $X$. Put
$E=\gen{ab,bc}$ and $M=\gen{a,b,c}$. If a nonempty $E$-invariant set
$T$ satisfies $2|T|<|X|$, then $M$ is intransitive on $X$.
\end{lemma}
\begin{proof}
Conjugation by $b$ inverts each of $ab$ and $bc$; hence $b$ normalizes $E$.
Furthermore $a=(ab)b$ and $c=b(bc)$, so $M=\gen{E,b}$. For $g\in E$,
\[
 (T^b)^g=T^{bg}=T^{(bgb)b}=(T^{bgb})^b=T^b.
\]
Thus both $T$ and $T^b$ are $E$-invariant. Since $b^2=1$, $b$ interchanges
them. Consequently $U=T\cup T^b$ is $M$-invariant and
\[
 0<|T|\leq |U|\leq2|T|<|X|.
\]
It is a nonempty proper invariant set, which proves the assertion.
No simplicity assumption is used.
\end{proof}

\begin{proposition}
Under the exact imported original-group and orbit premises, the supplied
strict certificate for free witness $290$, owner $579$, with row digest
\texttt{a5e7343cb4c3ca43d652371753233cd17df6b941251c8f5ec2a3791969c94ffc},
proves nongeneration provided its reviewed strict result certifies the
stated $351$-point invariant subset.
\end{proposition}
\begin{proof}
The strict checker reconstructs the same actual middle permutation from the
accepted equivariant procedural map; no uncomputed native family is trusted.
The Root case import supplies the exact original row, endpoint/frame equations,
nonidentity involution and inner-$G$ membership premises. The supplied-only
strict result is \texttt{PROPER} for this singleton row. It verifies closure
of the supplied set $T$ under the full positive $ab$ and $bc$ images.
Since those maps are bijections and $T$ is finite, these inclusions are
equalities. Thus $T$ is $E$-invariant. The independently reviewed projection
from the eight-row producer to the singleton strict carrier must bind the
same literal set, of cardinality $351$, rather than a newly substituted set.

The lemma yields an $M$-invariant union $U=T\cup T^b$ with
\[
 0<|U|\leq702<3510.
\]
The named group $G$ is transitive in this exact original action, by the
retained complete action bridge and maximal-subgroup argument. Hence
$M\ne G$. Imported membership gives $M\leq G$, so $M<G$.

The retained original $H$-orbit and two-fibre identification transfers the
result to the represented original scope. In particular, conjugation by
$H\leq A$ preserves properness in $G$ because $G\trianglelefteq A$.
The owner integer is an identifier, not an interpreted group word. No new
orbit, native map, recognition or transitivity calculation is used here.
\end{proof}

\paragraph{Attribution and limits.}
The invariant-set union and transitivity arguments are borrowed from
\emph{Nongeneration from small invariant subsets}, retained SHA256
\texttt{abc0f4affac1a88fbe836a4c54d6ad4d64328ed1d403c4487390e533f5baf742}.
That text uses $T\cup T^a$ and concerns earlier witnesses $249$ and $250$;
the displayed proof above explicitly supplies the equally valid $b$ version
and does not retag those old applications. The named maximal-subgroup
classification and its correction remain attributed inputs, as do the
original member, model, action, $H=B$ and two-fibre identifications.
This application concerns generation only. It supplies no intersection
property, Aut/reversal, class, census, manuscript, conjecture or Lean claim.
The seven other producer rows remain unpaid. Strict execution review,
Root mathematical adoption, logical peer review, retention and typed ledger
publication are separate requirements.
\end{document}
